\subsection{Low-Rank Approximation and Pseudoinverses}

Throughout this subsection, $\mathbb{F}\in\{\mathbb{R},\mathbb{C}\}$; every matrix dimension is positive.

\begin{proposition} \label{prop:rank-one-svd-orthogonality-and-truncation-error} \hlt{Rank-One SVD Terms and Truncation Error}\\
Let $A=X\Sigma_rY^*=\sum\limits_{j=1}^{r}\sigma_j\mathbf{x}_j\mathbf{y}_j^*$ be a compact SVD, with $\sigma_1\geq\cdots\geq\sigma_r>0$. For $k\in\{0,\ldots,r\}$, set $A_k\defeq\sum\limits_{j=1}^{k}\sigma_j\mathbf{x}_j\mathbf{y}_j^*$, where $A_0\defeq0$.
\begin{enumerate}[label=(\roman*)]
\setlength{\itemsep}{0pt}
\item $(\forall i,j\in J_r)\,\innerprod[F]{\sigma_i\mathbf{x}_i\mathbf{y}_i^*,\sigma_j\mathbf{x}_j\mathbf{y}_j^*}=\sigma_i^2\delta_{ij}$.
\item $\norm{A}_F^2=\sum\limits_{j=1}^{r}\sigma_j^2$ and $\norm{A-A_k}_F^2=\sum\limits_{j=k+1}^{r}\sigma_j^2$.
\item $k<r\Longrightarrow\norm{A-A_k}_2=\sigma_{k+1}$, while $k=r\Longrightarrow\norm{A-A_k}_2=0$.
\end{enumerate}
\begin{proof}
\begin{enumerate}[label=(\roman*)]
\setlength{\itemsep}{0pt}
\item $(\forall i,j\in J_r)\,\innerprod[F]{\sigma_i\mathbf{x}_i\mathbf{y}_i^*,\sigma_j\mathbf{x}_j\mathbf{y}_j^*}=\sigma_i\sigma_j(\mathbf{y}_i^*\mathbf{y}_j)(\mathbf{x}_j^*\mathbf{x}_i)=\sigma_i^2\delta_{ij}$ by orthonormality.
\item The first identity holds by Proposition \ref{prop:basic-singular-value-identities}(ii). For $k<r$, the residual is $A-A_k=\sum\limits_{j=k+1}^{r}\sigma_j\mathbf{x}_j\mathbf{y}_j^*$.\\
Its positive singular values are $\sigma_{k+1},\ldots,\sigma_r$ by Theorem \ref{thm:compact-singular-value-decomposition}(ii).\\
The second identity now holds by Proposition \ref{prop:basic-singular-value-identities}(ii); for $k=r$, it is $0=0$.
\item If $k<r$, then $\norm{A-A_k}_2=\sigma_{k+1}$ by (ii) and Theorem \ref{thm:spectral-norm-largest-singular-value}(i); if $k=r$, then $A-A_k=0$.
\end{enumerate}
\hlt{Structural Commentary}: Frobenius orthogonality separates rank-one terms; compact SVD identifies the residual spectrum.\\
\hlt{Key Insight}: Truncation leaves exactly the singular-value tail.
\end{proof}
\end{proposition}

Figure \ref{fig:truncated-svd-energy} separates the retained and discarded singular-value terms.

\begin{illustrationfigure}[H]
  \centering
  \begin{tikzpicture}[
  >=Stealth,
  every node/.style={font=\small},
  kept/.style={draw=blue!70!black, fill=blue!10, rounded corners=2pt, minimum width=0.65cm, minimum height=2.2cm},
  tail/.style={draw=red!70!black, fill=red!10, rounded corners=2pt, minimum width=0.65cm}
]
  \node at (-0.45,1.35) {$\sigma_j$};
  \draw[->, gray!70] (0,0) -- (8.4,0);
  \draw[->, gray!70] (0,0) -- (0,2.65);

  \foreach \x/\h/\lab in {0.7/2.25/1,1.55/1.8/2,2.4/1.45/k} {
    \node[kept, minimum height=\h cm] at (\x,\h/2) {};
    \node[blue!70!black] at (\x,-0.3) {$\lab$};
  }
  \foreach \x/\h/\lab in {3.45/1.05/k+1,4.3/0.82/,5.15/0.65/,6.0/0.52/r} {
    \node[tail, minimum height=\h cm] at (\x,\h/2) {};
    \node[red!70!black] at (\x,-0.3) {$\lab$};
  }

  \draw[dashed, gray!65] (2.9,0) -- (2.9,2.45);
\end{tikzpicture}

  \caption{Bars left of the dashed divider form $A_k$; the right-hand bars satisfy $\norm{A-A_k}_F^2=\sum\limits_{j=k+1}^{r}\sigma_j^2$.}
  \label{fig:truncated-svd-energy}
\end{illustrationfigure}

\begin{theorem} \label{thm:eckart-young-frobenius-spectral} \hlt{Eckart--Young Low-Rank Approximation}\\
Let $A\in\operatorname{Mat}(m,n,\mathbb{F})$ have rank $r$, set $q\defeq\min(m,n)$. Write $\sigma_1\geq\cdots\geq\sigma_q\geq0$ for its singular values, fix a full SVD $A=V\Sigma W^*$.\\
Write $\mathbf{v}_j,\mathbf{w}_j$ for the columns of $V,W$, and set $(\forall j\in J_n\setminus J_q)\,\sigma_j\defeq0$. For $0\leq k<q$, set $A_k\defeq\sum\limits_{j=1}^{k}\sigma_j\mathbf{v}_j\mathbf{w}_j^*$.\\
If $B\in\operatorname{Mat}(m,n,\mathbb{F})$ and $\operatorname{rank}(B)\leq k$, then:
\begin{enumerate}[label=(\roman*)]
\setlength{\itemsep}{0pt}
\item $\norm{A-A_k}_2\leq\norm{A-B}_2$;
\item $\norm{A-A_k}_F\leq\norm{A-B}_F$.
\end{enumerate}
\begin{proof}
If $k\geq r$, then $A_k=A$ and $\norm{A-A_k}_2=\norm{A-A_k}_F=0$, so both claims hold.\\
For $k<r$, the first $r$ columns of $V,W$ form a compact SVD of $A$ by Theorem \ref{thm:compact-singular-value-decomposition}(i). Its rank-$k$ truncation is $A_k$.
\begin{enumerate}[label=(\roman*)]
\setlength{\itemsep}{0pt}
\item $\dim\ker(B)\geq n-k$ by Corollary \ref{cor:matrix-rank-nullity}(i). $\ker(B)\cap\operatorname{span}(\mathbf{w}_1,\ldots,\mathbf{w}_{k+1})\neq\{\mathbf{0}\}$ by Corollary \ref{cor:large-subspaces-intersect}(i).\\
Fix a unit vector $\mathbf{x}=\sum\limits_{i=1}^{k+1}c_i\mathbf{w}_i$ in this intersection. Then $\sum\limits_{i=1}^{k+1}\abs{c_i}^2=1$ and $\norm{A\mathbf{x}}_2^2=\sum\limits_{i=1}^{k+1}\sigma_i^2\abs{c_i}^2\geq\sigma_{k+1}^2$ by orthonormality.\\
$\norm{A-B}_2\geq\norm{(A-B)\mathbf{x}}_2=\norm{A\mathbf{x}}_2\geq\sigma_{k+1}$ by $B\mathbf{x}=\mathbf{0}$, spectral-norm definition. $\sigma_{k+1}=\norm{A-A_k}_2$ by Proposition \ref{prop:rank-one-svd-orthogonality-and-truncation-error}(iii).
\item $\dim\ker(B)\geq n-k$ by Corollary \ref{cor:matrix-rank-nullity}(i); choose $\mathcal{U}\subseteq\ker(B)$ with $\dim\mathcal{U}=n-k$.\\
An orthonormal basis $\mathbf{z}_1,\ldots,\mathbf{z}_{n-k}$ of $\mathcal{U}$ exists by Corollary \ref{cor:finite-dimensional-inner-product-space-has-orthonormal-basis}. Extend it to $\mathbf{z}_1,\ldots,\mathbf{z}_n$ by Corollary \ref{cor:orthonormal-list-extends-to-basis}.\\
For $Z\defeq[\mathbf{z}_1\ \cdots\ \mathbf{z}_n]$ and $P\defeq P_{\mathcal{U}}$, $Z$ is unitary by Theorem \ref{thm:unitary-matrix-equivalent-conditions}(i)--(ii). $\norm{A-B}_F^2=\norm{(A-B)Z}_F^2$ by Proposition \ref{prop:unitary-invariance-of-singular-values}(iii).\\
$\norm{(A-B)Z}_F^2\geq\sum\limits_{j=1}^{n-k}\norm{A\mathbf{z}_j}_2^2$ by the column formula for $\norm{\cdot}_F$ and $(\forall j\in J_{n-k})\,B\mathbf{z}_j=\mathbf{0}$.\\
Set $(\alpha_i)_{i\in J_n}\defeq(\norm{P\mathbf{w}_i}_2^2)_{i\in J_n}$. Then $\sum\limits_{j=1}^{n-k}\norm{A\mathbf{z}_j}_2^2=\sum\limits_{i=1}^{n}\sigma_i^2\alpha_i$ by orthonormality and Proposition \ref{prop:orthogonal-projection-properties}(ix).\\
$(\forall i\in J_n)\,0\leq\alpha_i\leq1$ by Proposition \ref{prop:orthogonal-projection-properties}(viii). $\sum\limits_{i=1}^{n}\alpha_i=n-k$ by Theorem \ref{thm:parseval-identity-finite-dimensional}. Thus $\beta\defeq\sum\limits_{i=1}^{k}\alpha_i=\sum\limits_{i=k+1}^{n}(1-\alpha_i)$.\\
$\sum\limits_{i=1}^{n}\sigma_i^2\alpha_i-\sum\limits_{i=k+1}^{n}\sigma_i^2\geq\sigma_{k+1}^2(\beta-\beta)=0$ by singular-value ordering. $\norm{A-B}_F^2\geq\sum\limits_{i=k+1}^{n}\sigma_i^2=\norm{A-A_k}_F^2$ by Proposition \ref{prop:rank-one-svd-orthogonality-and-truncation-error}(ii).
\end{enumerate}
\hlt{Structural Commentary}: Kernel intersection controls spectral loss; projection weights majorize the Frobenius tail.\\
\hlt{Key Insight}: Rank $k$ preserves at most $k$ singular directions.
\end{proof}
\end{theorem}

\begin{computationmethod} \label{meth:truncated-svd-approximation} \hlt{Truncated SVD Approximation}\\
Given $A\in\operatorname{Mat}(m,n,\mathbb{F})$ with compact SVD $A=X\Sigma_rY^*$ and $k\in\{0,\ldots,r-1\}$:
\begin{enumerate}[label=(\roman*)]
\setlength{\itemsep}{0pt}
\item set $\mathcal{R}_k\defeq\{B\in\operatorname{Mat}(m,n,\mathbb{F}):\operatorname{rank}(B)\leq k\}$;
\item form $A_k\defeq\sum\limits_{j=1}^{k}\sigma_j\mathbf{x}_j\mathbf{y}_j^*\in\mathcal{R}_k$;
\item compute $\norm{A-A_k}_F^2=\sum\limits_{j=k+1}^{r}\sigma_j^2$ and $\norm{A-A_k}_2=\sigma_{k+1}$;
\item certify $(\forall\nu\in\{2,F\})\,\norm{A-A_k}_{\nu}=\min\limits_{B\in\mathcal{R}_k}\norm{A-B}_{\nu}$ by (ii) and Theorem \ref{thm:eckart-young-frobenius-spectral}.
\end{enumerate}
\end{computationmethod}

The same compact SVD data defines a canonical inverse on the nonzero singular subspaces.

\begin{definition} \label{def:moore-penrose-pseudoinverse} \hlt{Moore--Penrose Pseudoinverse}\\
For $A\in\operatorname{Mat}(m,n,\mathbb{F})$, define $A^{\dagger}\in\operatorname{Mat}(n,m,\mathbb{F})$ by:
\begin{enumerate}[label=(\roman*)]
\setlength{\itemsep}{0pt}
\item $A=0\Longrightarrow A^{\dagger}\defeq0_{n,m}$;
\item $\operatorname{rank}(A)=r\geq1\Longrightarrow A^{\dagger}\defeq Y\Sigma_r^{-1}X^*$ for a compact SVD $A=X\Sigma_rY^*$.
\end{enumerate}
The matrix $A^{\dagger}$ is the \hlt{pseudoinverse} or \hlt{Moore--Penrose pseudoinverse} of $A$.
\end{definition}

\begin{proposition} \label{prop:pseudoinverse-well-defined} \hlt{Pseudoinverse is Well Defined}\\
The definition of $A^{\dagger}$ does not depend on the chosen compact SVD of $A$.
\begin{proof}
The assertion is immediate for $A=0$. Suppose $\operatorname{rank}(A)=r\geq1$. Let $A=X_1\Sigma_rY_1^*=X_2\Sigma_rY_2^*$ be compact SVDs.\\
Unitary $U$ satisfies $X_1=X_2U$, $Y_1=Y_2U$, $U\Sigma_r=\Sigma_rU$ by Theorem \ref{thm:compact-svd-unitary-ambiguity}. $U\Sigma_r^{-1}U^*=\Sigma_r^{-1}$; $Y_1\Sigma_r^{-1}X_1^*=Y_2U\Sigma_r^{-1}U^*X_2^*=Y_2\Sigma_r^{-1}X_2^*$.\\
\hlt{Structural Commentary}: Compact-SVD ambiguity commutes with reciprocal scaling and cancels between the singular bases.\\
\hlt{Key Insight}: Equal-value blocks also commute with $\Sigma_r^{-1}$.
\end{proof}
\end{proposition}

\begin{corollary} \label{cor:pseudoinverse-from-full-svd} \hlt{Pseudoinverse from a Full SVD}\\
Let $A=V\Sigma W^*$ be a full SVD of $A\in\operatorname{Mat}(m,n,\mathbb{F})$ with rank $r$, and set $D_r\defeq\{(j,j):j\in J_r\}$.\\
Define $\Sigma^{\dagger}\in\operatorname{Mat}(n,m,\mathbb{F})$ by $(\forall j\in J_r)\,(\Sigma^{\dagger})_{jj}\defeq\sigma_j^{-1}$. $(\forall(i,j)\in(J_n\times J_m)\setminus D_r)\,(\Sigma^{\dagger})_{ij}\defeq0$. Then $A^{\dagger}=W\Sigma^{\dagger}V^*$.
\begin{proof}
If $r=0$, then $A=0$, $\Sigma^{\dagger}=0_{n,m}$, and $W\Sigma^{\dagger}V^*=A^{\dagger}=0_{n,m}$. Suppose $r\geq1$, and write $V=[X\ X_0]$ and $W=[Y\ Y_0]$.\\
Then $A=X\Sigma_rY^*$ and $W\Sigma^{\dagger}V^*=Y\Sigma_r^{-1}X^*$ by block multiplication. Hence $A^{\dagger}=W\Sigma^{\dagger}V^*$ by the defining compact-SVD formula.\\
\hlt{Structural Commentary}: Full-SVD block multiplication isolates the compact factors and their reciprocal scaling.\\
\hlt{Key Insight}: Only positive diagonal entries are inverted.
\end{proof}
\end{corollary}

\begin{theorem} \label{thm:moore-penrose-basic-properties} \hlt{Basic Properties of the Pseudoinverse}\\
Let $A\in\operatorname{Mat}(m,n,\mathbb{F})$. Then:
\begin{enumerate}[label=(\roman*)]
\setlength{\itemsep}{0pt}
\item $\operatorname{col}(A^{\dagger})=\operatorname{col}(A^*)$ and $\ker(A^{\dagger})=\ker(A^*)$;
\item $AA^{\dagger}$ is the orthogonal projection of $\mathbb{F}^m$ onto $\operatorname{col}(A)$;
\item $A^{\dagger}A$ is the orthogonal projection of $\mathbb{F}^n$ onto $\operatorname{col}(A^*)$;
\item $AA^{\dagger}A=A$ and $A^{\dagger}AA^{\dagger}=A^{\dagger}$;
\item $(A^{\dagger})^{\dagger}=A$;
\item if $m=n$ and $A$ is invertible, then $A^{\dagger}=A^{-1}$.
\end{enumerate}
\begin{proof}
If $A=0$, then (i)--(v) hold directly and the hypothesis of (vi) fails.\\
Suppose $\operatorname{rank}(A)=r\geq1$. Fix a compact SVD $A=X\Sigma_rY^*$; then $A^{\dagger}=Y\Sigma_r^{-1}X^*$ by the defining formula.
\begin{enumerate}[label=(\roman*)]
\setlength{\itemsep}{0pt}
\item $\operatorname{col}(A^{\dagger})=\operatorname{col}(Y)$ by $A^{\dagger}=Y\Sigma_r^{-1}X^*$ and $A^{\dagger}X\Sigma_r=Y$. $\ker(A^{\dagger})=\ker(X^*)$ by $Y^*A^{\dagger}=\Sigma_r^{-1}X^*$ and invertibility of $\Sigma_r$.\\
$A^*=Y\Sigma_rX^*$ is a compact SVD. $\operatorname{col}(A^*)=\operatorname{col}(Y)$ and $\ker(A^*)=\ker(X^*)$ by Theorem \ref{thm:compact-svd-subspaces}(i)--(ii).
\item $AA^{\dagger}=XX^*=P_{\operatorname{col}(X)}=P_{\operatorname{col}(A)}$ by Proposition \ref{prop:orthogonal-projection-properties}(ix) and Theorem \ref{thm:compact-svd-subspaces}(i).
\item $A^{\dagger}A=YY^*=P_{\operatorname{col}(Y)}=P_{\operatorname{col}(A^*)}$ by Proposition \ref{prop:orthogonal-projection-properties}(ix) and Theorem \ref{thm:compact-svd-subspaces}(i) for $A^*=Y\Sigma_rX^*$.
\item $AA^{\dagger}A=XX^*X\Sigma_rY^*=A$ and $A^{\dagger}AA^{\dagger}=YY^*Y\Sigma_r^{-1}X^*=A^{\dagger}$ by $X^*X=Y^*Y=I_r$.
\item Let $R\in\operatorname{Mat}(r,r,\mathbb{F})$ be the reversal permutation matrix, so $(\forall j\in J_r)\,R\mathbf{e}_j=\mathbf{e}_{r+1-j}$.\\
Then $A^{\dagger}=(YR)(R^*\Sigma_r^{-1}R)(XR)^*$ is a compact SVD with decreasing positive diagonal.\\
Hence $(A^{\dagger})^{\dagger}=(XR)(R^*\Sigma_rR)(YR)^*=X\Sigma_rY^*=A$ by the defining compact-SVD formula.
\item If $A$ is square and invertible, then $A^*$ is invertible by Proposition \ref{prop:adjoints-inverses-and-polynomials}(i).\\
$\operatorname{col}(A)=\operatorname{col}(A^*)=\mathbb{F}^n$ by Corollary \ref{cor:square-matrix-invertible-iff-full-rank}. $AA^{\dagger}=A^{\dagger}A=I_n$ by (ii)--(iii), so $A^{\dagger}=A^{-1}$.
\end{enumerate}
\hlt{Structural Commentary}: Compact-SVD multiplication identifies the two projectors; reversed reciprocal scaling recovers $A$.\\
\hlt{Key Insight}: Positive singular coordinates invert; null coordinates remain zero.
\end{proof}
\end{theorem}

Figure \ref{fig:pseudoinverse-four-spaces} identifies the active inverse pair and the annihilated kernels.

\begin{illustrationfigure}[H]
  \centering
  \begin{tikzpicture}[
  >=Stealth,
  every node/.style={font=\small},
  space/.style={draw=gray!70, fill=gray!8, rounded corners=3pt, minimum width=1.75cm, minimum height=0.7cm, align=center},
  map/.style={->, line width=1pt, blue!70!black}
]
  \node[space] (kerA) at (0,0) {$\ker(A)$};
  \node[space] (row) at (2.45,0) {$\operatorname{col}(A^*)$};
  \node[space] (col) at (4.95,0) {$\operatorname{col}(A)$};
  \node[space] (kerAst) at (7.4,0) {$\ker(A^*)$};
  \node[red!70!black] at (0,-0.65) {$A=0$};
  \node[red!70!black] at (7.4,-0.65) {$A^{\dagger}=0$};

  \draw[map] (row) -- node[above] {$A$} (col);
  \draw[map] (col) -- node[below] {$A^{\dagger}$} (row);
\end{tikzpicture}

  \caption{The arrows invert the active singular subspaces; the outer boxes mark $\ker(A)$ and $\ker(A^*)$.}
  \label{fig:pseudoinverse-four-spaces}
\end{illustrationfigure}

\begin{theorem} \label{thm:penrose-equations-characterize-pseudoinverse} \hlt{Penrose Characterization of the Pseudoinverse}\\
Let $A\in\operatorname{Mat}(m,n,\mathbb{F})$, $B\in\operatorname{Mat}(n,m,\mathbb{F})$. Set $\mathsf{P}_A(B)\defeq[ABA=A\land BAB=B\land(AB)^*=AB\land(BA)^*=BA]$; $B=A^{\dagger}\Longleftrightarrow\mathsf{P}_A(B)$.
\begin{proof}
\begin{enumerate}[label=\roman*.]
\setlength{\itemsep}{0pt}
\item[$(\Rightarrow)$] Suppose $B=A^{\dagger}$. The first two conjuncts hold by Theorem \ref{thm:moore-penrose-basic-properties}(iv).\\
Matrices $AB$ and $BA$ are orthogonal projections by Theorem \ref{thm:moore-penrose-basic-properties}(ii)--(iii). $(AB)^*=AB$ and $(BA)^*=BA$ by Theorem \ref{thm:orthogonal-projection-self-adjoint-idempotent}.
\item[$(\Leftarrow)$] Suppose $\mathsf{P}_A(B)$, and set $\mathcal{C}\defeq\operatorname{col}(A)$ and $\mathcal{D}\defeq\operatorname{col}(A^*)$. Then $(AB)^2=ABAB=AB$ and $(AB)^*=AB$.\\
The relations $\operatorname{range}(AB)\subseteq\mathcal{C}$ and $(\forall\mathbf{x}\in\mathbb{F}^n)\,AB(A\mathbf{x})=A\mathbf{x}$ establish $\operatorname{range}(AB)=\mathcal{C}$.\\
$AB=P_{\mathcal{C}}$ by Theorem \ref{thm:orthogonal-projection-self-adjoint-idempotent}. For $BA$, $(BA)^2=BABA=BA$ and $(BA)^*=BA$.\\
$\ker(BA)\subseteq\ker(A)$ by $(\forall\mathbf{x}\in\mathbb{F}^n)\,[BA\mathbf{x}=\mathbf{0}\Rightarrow A\mathbf{x}=ABA\mathbf{x}=\mathbf{0}]$.\\
$\ker(A)\subseteq\ker(BA)$ by $(\forall\mathbf{x}\in\mathbb{F}^n)\,[A\mathbf{x}=\mathbf{0}\Rightarrow BA\mathbf{x}=\mathbf{0}]$; hence $\ker(BA)=\ker(A)$.\\
$\operatorname{range}(BA)=(\ker(BA))^{\perp}=(\ker A)^{\perp}=\mathcal{D}$ by Proposition \ref{prop:adjoint-kernel-range-orthogonality}(iii)--(iv). $BA=P_{\mathcal{D}}$ by Theorem \ref{thm:orthogonal-projection-self-adjoint-idempotent}.\\
The identities $AB=AA^{\dagger}$ and $BA=A^{\dagger}A$ hold by Theorem \ref{thm:moore-penrose-basic-properties}(ii)--(iii).\\
$B=BAB=B(AA^{\dagger})=(BA)A^{\dagger}=(A^{\dagger}A)A^{\dagger}=A^{\dagger}$ by $\mathsf{P}_A(B)$ and Theorem \ref{thm:moore-penrose-basic-properties}(iv).
\end{enumerate}
\hlt{Structural Commentary}: Self-adjoint idempotence identifies both range projections; their product identities collapse $B$.\\
\hlt{Key Insight}: $AB=AA^{\dagger}$ and $BA=A^{\dagger}A$ force $B=A^{\dagger}$.
\end{proof}
\end{theorem}

\begin{theorem} \label{thm:pseudoinverse-minimum-norm-least-squares} \hlt{Pseudoinverse and Minimum-Norm Least Squares}\\
For $A\in\operatorname{Mat}(m,n,\mathbb{F})$ and $\mathbf{y}\in\mathbb{F}^m$, set $\mathbf{x}_0\defeq A^{\dagger}\mathbf{y}$.
\begin{enumerate}[label=(\roman*)]
\setlength{\itemsep}{0pt}
\item $\mathbf{x}_0\in\operatorname{LS}(A,\mathbf{y})$ and $\operatorname*{arg\,min}\limits_{\mathbf{x}\in\operatorname{LS}(A,\mathbf{y})}\norm{\mathbf{x}}_2=\{\mathbf{x}_0\}$.
\item If $\mathbf{y}\in\operatorname{col}(A)$, then $A\mathbf{x}_0=\mathbf{y}$ and $\operatorname*{arg\,min}\limits_{\mathbf{x}\in\mathbb{F}^n:\,A\mathbf{x}=\mathbf{y}}\norm{\mathbf{x}}_2=\{\mathbf{x}_0\}$.
\end{enumerate}
\begin{proof}
Set $P\defeq AA^{\dagger}=P_{\operatorname{col}(A)}$ by Theorem \ref{thm:moore-penrose-basic-properties}(ii); then $A\mathbf{x}_0=P\mathbf{y}$.\\
$\mathbf{x}_0\in\operatorname{col}(A^{\dagger})=\operatorname{col}(A^*)=(\ker A)^{\perp}$ by Theorem \ref{thm:moore-penrose-basic-properties}(i) and Proposition \ref{prop:adjoint-kernel-range-orthogonality}(iii).
\begin{enumerate}[label=(\roman*)]
\setlength{\itemsep}{0pt}
\item $\mathbf{x}_0\in\operatorname{LS}(A,\mathbf{y})$ by $A\mathbf{x}_0=P\mathbf{y}$ and Theorem \ref{thm:orthogonal-projection-minimizes-distance}. $P_{(\ker A)^{\perp}}\mathbf{x}_0=\mathbf{x}_0$ by Proposition \ref{prop:orthogonal-projection-properties}(ii).\\
The minimum-norm identity holds by Proposition \ref{prop:least-squares-affine-set-and-minimum-norm}(iii).
\item If $\mathbf{y}\in\operatorname{col}(A)$, then $A\mathbf{x}_0=P\mathbf{y}=\mathbf{y}$ by Proposition \ref{prop:orthogonal-projection-properties}(ii).\\
Thus $\operatorname{LS}(A,\mathbf{y})=\{\mathbf{x}\in\mathbb{F}^n:A\mathbf{x}=\mathbf{y}\}$ by the zero minimum residual, and the claim holds by (i).
\end{enumerate}
\hlt{Structural Commentary}: $AA^{\dagger}$ fits the data; $\operatorname{col}(A^{\dagger})=(\ker A)^{\perp}$ selects the shortest preimage.\\
\hlt{Key Insight}: The pseudoinverse removes every kernel component.
\end{proof}
\end{theorem}

\begin{computationmethod} \label{meth:pseudoinverse-solution-method} \hlt{Minimum-Norm Least Squares by Pseudoinverse}\\
Given $A\in\operatorname{Mat}(m,n,\mathbb{F})$ and $\mathbf{y}\in\mathbb{F}^m$, compute the minimum-norm least-squares solution as follows:
\begin{enumerate}[label=(\roman*)]
\setlength{\itemsep}{0pt}
\item if $A=0$, set $A^{\dagger}=0_{n,m}$;
\item if $\operatorname{rank}(A)=r\geq1$, compute a compact SVD $A=X\Sigma_rY^*$ and form $A^{\dagger}=Y\Sigma_r^{-1}X^*$;
\item set $\mathbf{x}_0\defeq A^{\dagger}\mathbf{y}$;
\item certify $\operatorname*{arg\,min}\limits_{\mathbf{x}\in\operatorname{LS}(A,\mathbf{y})}\norm{\mathbf{x}}_2=\{\mathbf{x}_0\}$ by Theorem \ref{thm:pseudoinverse-minimum-norm-least-squares}(i);
\item if $\mathbf{y}\in\operatorname{col}(A)$, report $A\mathbf{x}_0=\mathbf{y}$ by Theorem \ref{thm:pseudoinverse-minimum-norm-least-squares}(ii).
\end{enumerate}
\end{computationmethod}
